\section{Many constraints in a three-dimensional matrix span}
\label{app:three-span}

The number of generators matters when the homogeneous matrices span a
three-dimensional space but more than three constraints are present.
This appendix gives an elementary directional refinement of the known
facet bound and an obstruction to extending the two-aggregation case to
arbitrary polyhedral coefficient cones. These are supporting consequences
and examples, rather than a new topological aggregation theorem.

Let $n\geq3$, $L=\spanop\{Q_1,\ldots,Q_m\}$ have dimension three,
and suppose that $L$ contains a positive definite matrix. Assume
$S\ne\varnothing$ and $\conv S\ne\R^n$. Define the matrix cone
\[
 \mathcal C_L=\cone\{Q_1,\ldots,Q_m\}\subset L,
 \qquad \mathcal P_L=L\cap\Sym^{n+1}_+.
\]
We distinguish $\mathcal C_L$ from the multiplier cone $\mathcal K$
in \eqref{eq:certificate-cone}. For $M\in L$, write
$q_M(x)=(x,1)\transpose M(x,1)$.
Evaluation at a fixed feasible lift is negative on every nonzero
nonnegative combination of the generators. It follows that $\mathcal C_L$
is pointed. A section on which the negative of that evaluation equals one
is the convex hull of the normalized generators, hence is a compact
polygon. If $\mathcal C_L$ has $k$ extreme rays, it consequently has
$k$ facets, each generated by two extreme rays. We may retain one
original matrix on each extreme ray: every omitted row is a nonzero
nonnegative combination of the retained rows, so the strict set and the
available aggregation matrices are preserved.

Blekherman, Dey, and Sun give the resulting $2k$ bound
\citep[Proposition~2.22]{BDSpreprint}. Their proof uses the facets of
this three-dimensional cone. The next proposition records how the count
depends on the chosen positive definite direction; its three-generator
antecedent is \citet[Proposition~8.6]{BDpreprint}.

\begin{proposition}[Directional facet bound]\label{prop:facet-bound}
Write an irredundant facet representation
\[
 \mathcal C_L=\{M\in L:\ell_j(M)\geq0,\ j=1,\ldots,k\}.
\]
For every $D\in L$ with $D\succ0$, put
$J_-(D)=\{j:\ell_j(D)<0\}$. Then $\conv S$ is described by
at most $2|J_-(D)|$ good aggregations. In particular,
\begin{equation}\label{eq:conditional-facet-bound}
 \mathcal P_L\not\subseteq-\mathcal C_L
 \quad\Longrightarrow\quad
 \text{at most }2k-2\text{ good aggregations suffice}.
\end{equation}
\end{proposition}
\begin{proof}
First, the full system has HHC. Choose three basis matrices for $L$.
On each homogeneous hyperplane their restrictions still admit a positive
definite real combination. That hyperplane has dimension $n\geq3$, so
the three-form convexity theorem \citep[Theorem~2.1]{Polyak1998}
gives a convex joint image. The image of the original $m$ forms is a
linear image of this convex image. The full good-aggregation description
therefore applies \citep[Theorem~2.9]{BDSpreprint}.

Since $D\succ0$ has positive evaluation at a feasible lift, it cannot
belong to $\mathcal C_L$. Thus $J_-(D)$ is nonempty. Let $M\in\mathcal C_L$
be a good aggregation matrix whose strict sublevel is not all of
$\R^n$. It has exactly one negative eigenvalue. Set
\[
 t_* =\min_{j\in J_-(D)}
       \frac{\ell_j(M)}{-\ell_j(D)},\qquad M_*=M+t_*D.
\]
Every facet inequality remains satisfied up to $t_*$, and at least one
indexed by $J_-(D)$ becomes an equality. Hence $M_*\in\mathcal C_L$
lies on a selected facet. Also $M_*\ne0$: for $t_*=0$ this follows
from $M\ne0$, and for $t_*>0$ the alternative would make
$M=-t_*D$ negative definite, contrary to its inertia.

The positive semidefinite improvement lemma
\citep[Proposition~9.1]{BDSpreprint} shows that $M_*$ is good and
$\{q_{M_*}<0\}\subseteq\{q_M<0\}$. To check its coefficient
hypothesis, choose nonnegative representations $M=Q_\lambda$ and
$M_*=Q_\mu$. Then $Q_{\mu-\lambda}=t_*D\succeq0$ and
$\lambda+(\mu-\lambda)=\mu\geq0$, with $\mu\ne0$. Thus every
nonredundant good cut is dominated by a good cut on a selected facet.
The intersection of all good cuts on these facets equals the hull:
it contains the hull by goodness and is contained in every original good
cut by this domination.

On each facet the available matrices are nonnegative combinations of
its two extreme generators. The pair-support reduction
\citep[Proposition~9.6]{BDSpreprint} describes the intersection of its
good cuts using at most two such cuts. Goodness here is relative to the
full set $S$, as in that proposition, not to the two-row subsystem.
Facets without good cuts contribute no inequality. Summing over the
selected facets proves $2|J_-(D)|$.

Finally, if $E\in\mathcal P_L\setminus(-\mathcal C_L)$, some
$\ell_j(E)>0$. Choose $D_0\succ0$ in $L$. For sufficiently small
$\delta>0$, $D=E+\delta D_0\succ0$ and $\ell_j(D)>0$.
Thus $|J_-(D)|\leq k-1$, proving
\eqref{eq:conditional-facet-bound}.
\end{proof}

Minimizing $|J_-(D)|$ over positive definite $D\in L$ can strengthen
the displayed bound for a particular cone. This is an existence bound;
no procedure for finding a best direction or the endpoint multipliers is
asserted. The proof uses strict aggregation directly and requires no
regularity or boundedness assumption on the original weak feasible set.

\begin{example}[An arbitrarily large necessary family of ellipsoids]
\label{ex:many-ellipsoids}
Let $m\geq3$, $n\geq2$, and $0<a_1<\cdots<a_m<1$. Put
$\rho=\sum_{j=2}^n x_j^2$ and define
\begin{equation}\label{eq:many-ellipsoids}
 f_i(x)=a_i x_1^2+(1-a_i^2)\rho-1,\qquad
 U=\{x:f_i(x)<0\text{ for all }i\}.
\end{equation}
Then $U$ is nonempty, open, bounded, and convex, since its defining
quadratic parts are positive definite and the origin is strictly feasible.
In the basis $(x_1^2,\rho,t^2)$ of the homogeneous span, the matrices
have coordinate vectors $(a_i,1-a_i^2,-1)$. For $a<b<c$,
\begin{equation}\label{eq:span-vandermonde}
 \det\begin{pmatrix}a&1-a^2&-1\\b&1-b^2&-1\\c&1-c^2&-1\end{pmatrix}
 =(b-a)(c-a)(c-b)>0.
\end{equation}
This follows by the elementary column operations that turn the matrix
into a Vandermonde matrix. Hence the span is exactly three-dimensional
and contains the positive definite form $x_1^2+\rho+t^2$.

For each $i$, choose $p_i$ with
\begin{equation}\label{eq:many-witness}
 (p_i)_1^2=\frac{2a_i}{1+a_i^2},\qquad
 \sum_{j=2}^n(p_i)_j^2=\frac1{1+a_i^2}.
\end{equation}
For example, take the positive square roots in coordinates one and two
and zeros elsewhere. Direct substitution gives the identity
\begin{equation}\label{eq:many-witness-slack}
 f_j(p_i)=-\frac{(a_j-a_i)^2}{1+a_i^2}.
\end{equation}
Thus, at $p_i$, a nonzero nonnegative aggregation is nonnegative
exactly when its multiplier is a positive multiple of $e_i$.
Since $p_i\notin U=\conv U$, every exact strict aggregation family,
even an infinite one, must contain that ray. The original $m$ rows are
good, so the minimum finite cardinality is exactly $m$. Each row has
exactly one negative homogeneous eigenvalue. The witness also exposes
its matrix ray in $\mathcal C_L$: all other generators have strictly
negative evaluation at $p_i$. In particular $k=m$.

The same necessity holds for finite nonstrict aggregation descriptions
of $\cl U$. Indeed, if a finite family omits ray $i$, all its
inequalities are strict at $p_i$ and remain satisfied in a neighborhood.
For $t>1$ close to one, $f_i(tp_i)=t^2-1>0$, so that neighborhood
contains a point outside $\cl U$. Conversely the original $m$ weak
inequalities describe $\cl U$: shrinking any weak feasible point
toward the origin makes every inequality strict. Thus $m$ is also the
minimum finite weak count. No indispensability assertion is made for
infinite weak families, where a missing ray can be approached by limits.
\end{example}

The complementary geometry $\mathcal P_L\subseteq-\mathcal C_L$ does
not imply a bound independent of $k$. In particular the two-aggregation
conclusion of \citet[Proposition~8.10]{BDpreprint}, proved for a
three-generator cone under its closed-set hypotheses, fails for general
polyhedral cones even when those geometric hypotheses hold.

\begin{proposition}[Obstruction in the complementary cone case]
\label{prop:complementary-many}
For every $m\geq3$ and $n\geq3$, there is a strictly feasible
system of $m+2$ quadratic inequalities whose three-dimensional
homogeneous span contains a
positive definite matrix, for which $\mathcal P_L\subseteq-\mathcal C_L$
and $k=m+2$, but every exact strict aggregation description of its hull
needs the $m=k-2$ specified rays. Exactly $m$ good aggregations suffice.
Its original weak feasible set is compact, has nonempty interior, equals
the closure of its interior, and has no points at infinity. Its minimum
finite weak aggregation count is also exactly $m$.
\end{proposition}
\begin{proof}
Start with \eqref{eq:many-ellipsoids} and append the two inequalities
\[
 g_1(x)=-x_1^2<0,\qquad g_2(x)=-\rho<0.
\]
The new strict set is
\[
 V=U\cap\{x:x_1\ne0,\ (x_2,\ldots,x_n)\ne0\}.
\]
The added rows change the strict set, but $\conv V=U$.
To see this, fix $x\in U$ and choose a ball about $x$ contained in
$U$. A sufficiently small vector $z$ can be chosen with
$z_1\notin\{x_1,-x_1\}$ and
$(z_2,\ldots,z_n)\notin\{(x_2,\ldots,x_n),-(x_2,\ldots,x_n)\}$.
Both $x+z$ and $x-z$ then belong to $V$. Thus $x$ is their midpoint.
The reverse inclusion follows from convexity of $U$. This also proves
nonemptiness and properness of the new hull.

In the coordinate basis $(x_1^2,\rho,t^2)$, let $E_1,E_2,E_3$
be the three basis matrices. The new matrix cone is
\[
 \mathcal C_L=\cone\{Q_1,\ldots,Q_m,-E_1,-E_2\}.
\]
At every $p_i$ from \eqref{eq:many-witness}, both added rows are
strictly negative, while \eqref{eq:many-witness-slack} is unchanged.
Each $Q_i$ therefore still spans an exposed extreme ray. The rays
$-E_1,-E_2$ are also extreme: every $Q_i$ has $t^2$ coefficient $-1$,
so a nonnegative representation of a matrix with zero $t^2$ coefficient
can use only $-E_1,-E_2$. These two rays are independent and extreme
in their two-dimensional cone. Since all generators have now been shown
extreme and distinct, $k=m+2$.

For every $i$ the exact identity
\begin{equation}\label{eq:negative-orthant-contained}
 Q_i+a_i(-E_1)+(1-a_i^2)(-E_2)=-E_3
\end{equation}
places $-E_3$ in $\mathcal C_L$. A matrix
$uE_1+vE_2+wE_3$ is positive semidefinite exactly when
$u,v,w\geq0$. Hence $-\mathcal P_L$ is the cone of
$-E_1,-E_2,-E_3$, and \eqref{eq:negative-orthant-contained} proves
$\mathcal P_L\subseteq-\mathcal C_L$.

Any nonzero nonnegative aggregation of the enlarged system that is
not supported only on original row $i$ is strictly negative at $p_i$.
That point is excluded from $\conv V=U$. Therefore every exact strict
aggregation description must contain each of the $m$ original rays,
even if additional infinitely many aggregations are allowed. Conversely,
the original $m$ convex inequalities are good for $V$ and describe
its hull. They attain the bound $m$.

The added weak inequalities hold globally. The original weak set is
therefore still $C=\{f_i\leq0\ \forall i\}=\cl U$.
It is compact, has interior $U$, and equals the closure of that interior.
For the interior equality, a point with $f_i=0$ lies on the boundary
of its positive definite ellipsoid and cannot be interior to $C$.
Each original leading matrix is positive definite, so a common
nonpositive direction at infinity must be zero. The same finite-family
neighborhood argument at $p_i$ proves that a weak aggregation description
of $C$ needs all $m$ original rays. They give the matching description.
\end{proof}

Thus the complementary hypothesis alone does not support a general-cone
two-bound, even for regular compact weak systems. This disproves that
possible shortcut to improving the facet count. It does not disprove an
unconditional $2k-2$ bound: the necessary count $k-2$ in this example
is smaller. Our established general upper bound remains the cited $2k$,
with \eqref{eq:conditional-facet-bound} and the directional count when
their hypotheses apply. No optimal count for arbitrary cones is asserted.
Nor can one obtain such a count simply by triangulating the coefficient
cone and intersecting subsystem hulls: convex hull does not commute with
intersection, and goodness is relative to the original full feasible set.
The inertia condition also fails to be closed under addition: the two
matrices $\operatorname{diag}(-4,1,1,0)$ and
$\operatorname{diag}(1,-4,1,0)$ each have one negative eigenvalue,
whereas their sum has two. Consequently their individual inertia bounds
cannot justify a reduction that adds good matrices arbitrarily.
